% =============================================================================
%  latex-agentic — single-page resume template
%  Engine: XeLaTeX ONLY (compile with `latexmk` or `make pdf`).
%
%  Design goals:
%    * One page, clear visual hierarchy, a single tasteful accent colour.
%    * ATS-friendly: real selectable text, standard section names, no tables
%      for the main content flow, no icon fonts used as content.
%    * Portable: built on the stock `article` class plus widely available
%      packages — NO exotic resume document classes. Works on any TeX Live.
%
%  HOW TO EDIT (read this!):
%    1. Replace the persona block under "PERSONAL DETAILS" with your own.
%    2. Each \role / \entry / bullet is plain text — just overwrite it.
%    3. Keep it to one page: trim bullets before shrinking fonts.
%    4. The accent colour is one line (\definecolor{accent}); change it freely.
% =============================================================================
\documentclass[a4paper,11pt]{article}

% --- Geometry: tight, even margins for a dense but breathable single page ---
\usepackage[
  top=1.4cm, bottom=1.4cm, left=1.6cm, right=1.6cm
]{geometry}

% --- Fonts (XeLaTeX / fontspec) ---------------------------------------------
% A sans-serif face reads as modern and stays crisp on screen and in print.
% We try several fonts in order and use the first one that exists, so this
% file compiles whether or not you have the "nice" fonts installed:
%
%   1. Source Sans 3 / Source Sans Pro  (clean humanist sans, may need install)
%   2. Lato                             (ships with TeX Live: tlmgr "lato")
%   3. TeX Gyre Heros                   (ships with TeX Live; Helvetica-like)
%   4. Latin Modern Sans                (always present — last resort)
%
% \IfFontExistsTF tests by font *name* first; for the TeX Live faces we fall
% back to testing by *filename*, which always resolves inside the TeX tree.
\usepackage{fontspec}
\defaultfontfeatures{Ligatures=TeX}

\IfFontExistsTF{Source Sans 3}{%
  \setmainfont{Source Sans 3}%
}{\IfFontExistsTF{Source Sans Pro}{%
  \setmainfont{Source Sans Pro}%
}{\IfFontExistsTF{Lato}{%
  \setmainfont{Lato}%
}{\IfFontExistsTF{texgyreheros-regular.otf}{%
  \setmainfont{TeX Gyre Heros}[%
    Path=, % resolved via the TeX tree
    Extension=.otf,
    UprightFont=texgyreheros-regular,
    BoldFont=texgyreheros-bold,
    ItalicFont=texgyreheros-italic,
    BoldItalicFont=texgyreheros-bolditalic]%
}{%
  \setmainfont{Latin Modern Sans}% always available
}}}}

% --- Colour, sectioning, lists, links ---------------------------------------
\usepackage{xcolor}
% Accent colour — one tasteful blue-teal. Change these RGB values to rebrand.
\definecolor{accent}{HTML}{1F6F8B}
\definecolor{rulegrey}{HTML}{D0D5DA}
\definecolor{textgrey}{HTML}{444B52}

\usepackage{titlesec}   % custom, compact section headings
\usepackage{enumitem}   % tight, configurable lists
\usepackage{microtype}  % subtle spacing polish
\usepackage[autostyle=true]{csquotes}

% hyperref LAST. colorlinks keeps the PDF clean (no boxes) and ATS-safe.
\usepackage[
  unicode=true,
  hidelinks,            % no coloured boxes; links still clickable
  pdftitle={Jane Doe — Software Engineer},
  pdfauthor={Jane Doe}
]{hyperref}

% --- Section heading style: accent text + a thin full-width rule ------------
% Note: the heading text is upper-cased via the 4th argument (\MakeUppercase
% wraps only the title), NOT via a prefix in the format spec — otherwise the
% colour names in the trailing rule would be upper-cased too and break.
\titleformat{\section}
  {\color{accent}\large\bfseries}             % look (no uppercase here)
  {}{0pt}                                      % no number, no label sep
  {\MakeUppercase}                             % apply uppercase to title only
  [\vspace{-0.55em}{\color{rulegrey}\titlerule[1pt]}] % rule under heading
\titlespacing*{\section}{0pt}{0.9em}{0.5em}

% --- Reusable building blocks (edit content, not these macros) --------------
% \entry{Title}{Right-hand note (e.g. dates)}{Subtitle / employer}{Location}
\newcommand{\entry}[4]{%
  \par\smallskip
  \noindent\textbf{#1}\hfill{\color{textgrey}\small #2}\\[1pt]
  {\color{accent}\small #3}\hfill{\color{textgrey}\small\itshape #4}\par
}

% Compact bullet list used under each entry.
\newlist{points}{itemize}{1}
\setlist[points]{
  leftmargin=1.2em, topsep=2pt, itemsep=1.5pt, parsep=0pt,
  label={\color{accent}\textbullet}
}

% A "Skills"-style labelled line: \skillline{Languages}{Go, Python, ...}
\newcommand{\skillline}[2]{%
  \noindent\textbf{#1:}~{#2}\par\smallskip%
}

\setlength{\parindent}{0pt}
\pagestyle{empty}  % no page numbers on a one-pager

\begin{document}

% =============================================================================
%  HEADER — name, role, and a single contact line with live links.
%  The contact line uses \href so email/GitHub/LinkedIn are clickable in the
%  PDF while still printing as plain, ATS-readable text.
% =============================================================================
% The new-resume skill replaces the sample name "Jane Doe" with the candidate's
% name (here and in the hyperref pdftitle/pdfauthor above). Editing by hand?
% Just overwrite it.
{\Huge\bfseries\color{accent} Jane Doe}\\[2pt]
{\large\color{textgrey} Senior Software Engineer}\\[6pt]
{\small
  \href{mailto:jane.doe@example.com}{jane.doe@example.com}
  \quad\textbar\quad +1 (555) 010-2920
  \quad\textbar\quad San Francisco, CA
  \\[2pt]
  \href{https://github.com/janedoe}{github.com/janedoe}
  \quad\textbar\quad
  \href{https://www.linkedin.com/in/janedoe}{linkedin.com/in/janedoe}
  \quad\textbar\quad
  \href{https://janedoe.dev}{janedoe.dev}%
}

% =============================================================================
\section{Summary}
% 2–3 lines. Lead with seniority + specialism + a headline result.
Backend-leaning full-stack engineer with 8 years building reliable,
high-throughput services. Ship-focused, pragmatic about trade-offs, and happy
mentoring. Reduced infrastructure spend by 30\% while improving p99 latency on
a platform serving 4M daily users.

% =============================================================================
\section{Experience}
% Reverse-chronological. Quantify outcomes ("by 30\%", "to 120ms") — numbers
% survive ATS parsing and catch a human reader's eye.

\entry{Senior Software Engineer}{2021 -- Present}{Northwind Systems}{San Francisco, CA}
\begin{points}
  \item Led the migration of a monolith to 12 Go microservices, cutting
        mean deploy time from 40 minutes to under 6 and halving on-call pages.
  \item Designed an event-driven billing pipeline (Kafka, PostgreSQL) handling
        2M events/day with 99.98\% delivery accuracy.
  \item Mentored 5 engineers; introduced a code-review rubric that dropped
        post-release defects by 22\% quarter over quarter.
\end{points}

\entry{Software Engineer}{2017 -- 2021}{Cobalt Labs}{Austin, TX}
\begin{points}
  \item Built the company's first public REST + GraphQL API, adopted by 80+
        external partners within a year.
  \item Cut median API latency from 310ms to 120ms by introducing read
        replicas and a Redis caching layer.
  \item Automated the release pipeline (CI/CD), enabling daily deploys and
        reducing manual release effort by roughly 15 hours per week.
\end{points}

% =============================================================================
\section{Education}
\entry{B.Sc. in Computer Science}{2013 -- 2017}{University of Example}{Boston, MA}
\begin{points}
  \item Graduated with honours (GPA 3.8/4.0); focus on distributed systems.
\end{points}

% =============================================================================
\section{Skills}
% Grouped, compact lines. Keep them scannable; lead with your strongest stack.
\skillline{Languages}{Go, Python, TypeScript, SQL, Bash}
\skillline{Backend}{gRPC, REST, GraphQL, Kafka, PostgreSQL, Redis}
\skillline{Infrastructure}{Docker, Kubernetes, Terraform, AWS, GitHub Actions}
\skillline{Practices}{TDD, observability (OpenTelemetry), code review, mentoring}

% =============================================================================
\section{Projects}
\entry{ledgerlite — open-source double-entry ledger}{2023}{github.com/janedoe/ledgerlite}{}
\begin{points}
  \item A 4k-line Go library for embedded accounting; 1.2k GitHub stars and
        used in production by three startups.
\end{points}

\entry{texwatch — incremental LaTeX rebuilder}{2022}{github.com/janedoe/texwatch}{}
\begin{points}
  \item A tiny file-watcher that rebuilds LaTeX projects on save; featured in a
        popular newsletter, 600+ stars.
\end{points}

\end{document}
